\chapter{Plan de gestión del trabajo}\label{cap:Planificacion}

Este capítulo describe cómo se organiza el desarrollo del proyecto: la metodología adoptada, los recursos empleados, la gestión de la configuración, la planificación temporal y la estimación de costes y riesgos.

\section{Metodología de desarrollo}

Dado que el proyecto combina investigación aplicada en la selección y entrenamiento de modelos de lenguaje con el desarrollo de un sistema de software completo, se opta por una \textbf{metodología ágil basada en Kanban}~\cite{Vijayasarathy16}. Este enfoque es especialmente apropiado cuando los requisitos evolucionan a medida que se van obteniendo los resultados.

Kanban no impone iteraciones de duración fija ni \emph{sprints} formales, sino que organiza el trabajo en un tablero con columnas que representan los posibles estados de cada tarea: \textbf{Backlog} (identificadas pero no iniciadas), \textbf{En progreso} (activas), \textbf{En revisión} (completadas y pendientes de verificación) y \textbf{Hecho} (finalizadas). Esto permite visualizar el flujo de trabajo de forma continua sin la rigidez de otras metodologías.

El límite de trabajo en curso (\emph{Work In Progress}, WIP) se fija en dos tareas simultáneas en la columna de \emph{En progreso}, una restricción que obliga a completar las tareas antes de comenzar nuevas y reduce el coste de los cambios de contexto.

\subsection{Flujo de trabajo en el tablero Kanban}

Cada tarea se representa como una tarjeta con descripción breve, componente al que pertenece y criterio de aceptación. Las tarjetas transitan por cuatro columnas:

\begin{itemize}
 \item \textbf{Backlog}: tareas identificadas y priorizadas pero no iniciadas. Al comienzo de cada sesión se selecciona la de mayor prioridad respetando el límite WIP.
 \item \textbf{En progreso}: tareas activas. El límite de dos simultáneas reduce la fragmentación del esfuerzo y facilita completar antes de comenzar.
 \item \textbf{En revisión}: implementación terminada, pendiente de verificación (pruebas, revisión del código o comprobación manual del comportamiento).
 \item \textbf{Hecho}: tareas completadas y verificadas. Una tarea no avanza hasta cumplir su criterio de aceptación.
\end{itemize}

La priorización del \emph{backlog} se revisa cada vez que se completa una tarea, no en ciclos fijos. Esto permite reaccionar con agilidad ante los resultados del entrenamiento: cuando los experimentos con un modelo base dan resultados insatisfactorios, la tarea de comparativa de modelos asciende en prioridad antes de que el desvío afecte al resto de la planificación.

\section{Recursos}

\subsection{Recursos humanos}

El proyecto se realiza íntegramente por un único desarrollador, con la dirección y supervisión de ambos tutores. El desarrollador asume todos los roles del proceso de desarrollo (analista, diseñador, desarrollador, responsable de pruebas y documentador).

\subsection{Recursos hardware}

\begin{itemize}
 \item \textbf{Equipos informáticos}: la máquina de desarrollo (Ryzen 9 7900X, RTX 4080) con Ubuntu 26.04 LTS y un portátil MacBook Pro M4, ambos utilizados para el desarrollo de los componentes del sistema (backend, frontend, scripts de datos, etc.) y para la ejecución de los distintos contenedores Docker.
 \item \textbf{GPU para entrenamiento}: El modelo BERT requiere para su entrenamiento el acceso a una GPU con soporte CUDA. Se utiliza la máquina de desarrollo mencionada. El entrenamiento se ejecuta en un entorno dockerizado sobre GPU local; la ejecución en la nube sería una alternativa que no se ha probado. En CPU, el entrenamiento tardaría entre 12 y 24 horas según el número de épocas (estimación, no medida); en GPU (RTX 4080), entre 1,7 y 134 minutos por ejecución (mediana de 24 minutos, sin contar el preentrenamiento), según \texttt{training\_runs.csv}.
 \item \textbf{Espacio en disco}: El corpus de CodiEsp y los artefactos del modelo entrenado requieren aproximadamente 5\,GB de almacenamiento; se dispondrá de al menos 200\,GB para garantizar espacio suficiente para logs, checkpoints y otros archivos temporales.
\end{itemize}

\subsection{Recursos software}

Las herramientas y tecnologías empleadas en el proyecto se agrupan en las siguientes categorías:

\textbf{Entorno de desarrollo:}
\begin{itemize}[noitemsep]
 \item Visual Studio Code y Spacemacs como editores principales.
 \item Docker para la ejecución de los contenedores en Ubuntu y macOS.
 \item Jupyter Notebook para el desarrollo interactivo de los experimentos de entrenamiento y análisis de datos.
 \item Make para la automatización de tareas comunes (entrenamiento, pruebas, generación de documentación).
 \item Chromium y Postman para la prueba manual de la API REST.
\end{itemize}

\textbf{Lenguajes y \emph{frameworks}:}
\begin{itemize}[noitemsep]
 \item Python 3.x con FastAPI, PyTorch y HuggingFace Transformers para el motor de inferencia.
 \item Elixir 1.19 con Phoenix 1.8 y Commanded para el backend.
 \item TypeScript con Next.js 16 y React 19 para el frontend.
 \item SQL (PostgreSQL 18) para la persistencia.
\end{itemize}

\textbf{Control de versiones y calidad:}
\begin{itemize}[noitemsep]
 \item Git y GitHub para el control de versiones y el alojamiento del repositorio.
 \item Husky y commitlint para la validación automática del formato de los mensajes de commit (Conventional Commits).
 \item GitHub Actions para la integración continua (CI) en las diferentes partes del proyecto.
\end{itemize}

\textbf{Pruebas:}
\begin{itemize}[noitemsep]
 \item ExUnit y excoveralls para las pruebas del backend en Elixir.
 \item Vitest para las pruebas del frontend en TypeScript.
 \item pytest para las pruebas del motor de IA y pytest con Playwright para las pruebas de extremo a extremo.
\end{itemize}

\textbf{Documentación:}
\begin{itemize}[noitemsep]
 \item \LaTeX{} para la redacción de la memoria del TFG.
 \item Markdown para la documentación técnica interna del repositorio.
\end{itemize}

\textbf{Infraestructura:}
\begin{itemize}[noitemsep]
 \item Docker Compose para la orquestación de los contenedores en desarrollo.
 \item Mailpit como servidor SMTP de captura de correo electrónico en el entorno local.
 \item GitHub Container Registry como caché de construcción de las imágenes en el CI.
 \item Docker Compose y Traefik (proxy inverso para la gestión de rutas) para el entorno de despliegue, publicado con un túnel de Cloudflare desde la propia estación de trabajo.
 \item Prometheus y Grafana para la monitorización del rendimiento y la salud del sistema desplegado.
 \item Sentry, un servicio externo (SaaS), para la gestión de errores y excepciones del sistema desplegado.
\end{itemize}

\section{Gestión de la configuración y aseguramiento de la calidad}

\subsection{Control de versiones}

El repositorio Git del proyecto~\cite{cie10repo2026} sigue el esquema de ramas habitual para este tipo de proyectos: una rama principal (\texttt{main}) que contiene el código con las pruebas del CI en verde, además de las ramas de trabajo (\texttt{feature branches}) para el desarrollo de cada funcionalidad, que se van integrando en la rama principal una vez están acabadas y validadas.

Al tratarse de un proyecto unipersonal, la integración de las ramas de trabajo se realiza mediante \textbf{fast-forward merge}, sin \emph{pull requests} ni \emph{merge commits}. Este flujo simplificado es apropiado para un único desarrollador por varias razones: el historial resultante es lineal y fácil de leer (\texttt{git log} muestra los cambios en orden cronológico sin bifurcaciones artificiales), se evita el ruido de los commits de fusión automáticos (\texttt{Merge branch 'feature/X' into main}) que no aportan información, y la revisión de código, necesaria en equipos para coordinar el trabajo, se realiza directamente durante el desarrollo. El uso de \emph{pull requests} en un proyecto unipersonal añadiría fricción sin beneficio real, ya que su función principal es facilitar la revisión entre pares, el control de acceso a la rama principal y la discusión asíncrona en equipos distribuidos.

Los mensajes de commit siguen el estándar \emph{Conventional Commits}~\cite{conventionalcommits2024}, validado automáticamente mediante el \emph{hook} \texttt{commit-msg} de Husky y commitlint. Este estándar impone un formato estructurado (\texttt{tipo(ámbito): descripción}) que facilita la lectura del historial y la generación automática de notas de versión.

\subsection{Integración continua}

Se configura un flujo de CI con GitHub Actions que se ejecuta en cada \emph{push} a la rama principal y en cada \emph{pull request}. El flujo incluye el análisis estático y el formato (ESLint, Credo, Dialyzer y Ruff), la compilación del frontend, las pruebas con cobertura del frontend, el backend y el motor de IA, las pruebas de extremo a extremo y la compilación de la memoria. La verificación del formato de los commits con commitlint se ejecuta en el CI solo en los \emph{pull requests}; en la rama principal la valida el \emph{hook} local de Husky.

\subsection{Aseguramiento de la calidad}

La calidad del código se procura mediante varios mecanismos complementarios:

\begin{itemize}
 \item \textbf{Pruebas automáticas}: La creación de la suite ExUnit para el backend (con informe de cobertura generado por excoveralls) y Vitest para el frontend.
 \item \textbf{Linting}: El uso y configuración de linters específicos para cada lenguaje (ESLint para TypeScript, Ruff para Python, Credo para Elixir) con reglas adaptadas al estilo del proyecto.
 \item \textbf{Revisión de código}: Al tratarse de un proyecto unipersonal, el código se revisa mediante la lectura sistemática antes de cada \emph{commit} significativo y mediante las comprobaciones automáticas del CI. La memoria, en cambio, se revisa además con una herramienta de apoyo basada en un asistente de IA (capítulo~\ref{cap:Conclusiones}).
 \item \textbf{Separación de entornos}: El uso de stubs y \emph{mocks} (objetos simulados) para el motor de inferencia (\texttt{ai\_engine\_mock}) permite desarrollar y evaluar el estado, tanto del backend como del frontend, sin necesidad del modelo BERT entrenado, evitando dependencias que dificulten las pruebas de integración.
 \item \textbf{Métricas de entrenamiento}: El proceso de entrenamiento registra las métricas de validación en cada época y genera curvas de aprendizaje, lo que permite detectar sobreajuste o problemas de convergencia de forma inmediata para actuar en consecuencia (ajuste de hiperparámetros, cambio de modelo base, etc.).
\end{itemize}

\section{Planificación}

El tablero Kanban se organiza en torno a seis \textbf{hitos o épicas} que agrupan el conjunto de tareas del proyecto y cada una representa un componente o fase del sistema. Las tareas se priorizan y completan de forma continua a lo largo del desarrollo.

\subsection{Épica 1: Adquisición y preparación de datos}

Tareas completadas:
\begin{itemize}[noitemsep]
 \item Descarga del corpus CodiEsp (diagnósticos, procedimientos y variante extendida).
 \item Desarrollo del script de \emph{scraping} sobre el portal eCIE-maps para obtener el catálogo CIE-10-ES.
 \item Análisis estadístico del corpus: Distribución de longitudes, frecuencia de códigos, desequilibrio de clases, palabras vacías y lematización.
 \item Análisis estadístico del catálogo CIE-10-ES: Cobertura, distribución por capítulos y uso de vocabulario médico específico.
 \item Generación de los ficheros CSV de entrenamiento, validación y prueba.
 \item Aumentación del corpus CodiEsp mediante retrotraducción con tres pivotes (EN, DE, FR) usando Azure Translator: cada pivote genera 500 paráfrasis adicionales que, sumadas a los 500 documentos originales, cuadruplican el corpus hasta 2.000 documentos (500 originales + 500 EN + 500 DE + 500 FR). La aumentación se evaluó y se descartó porque empeoró el resultado (Anexo~\ref{cap:AnexoL}).
\end{itemize}

\subsection{Épica 2: Modelo de clasificación}

Tareas completadas:
\begin{itemize}[noitemsep]
 \item Comparativa de los distintos modelos base apropiados para este uso.
 \item Implementación del clasificador plano multilabel basado en un modelo clínico preentrenado BERT.
 \item Configuración del pipeline de entrenamiento (función de pérdida, optimizador, \emph{scheduler}).
 \item Implementación del clasificador de diccionario como \emph{baseline}.
 \item Entrenamiento y selección del mejor \emph{checkpoint} por F1-micro en validación (por MAP en el modelo servido).
 \item Evaluación sobre el conjunto de prueba con métricas de precisión, recall y F1.
\end{itemize}

\subsection{Épica 3: Servicio de inferencia (AI Engine)}

Tareas completadas:
\begin{itemize}[noitemsep]
 \item Implementación del servicio FastAPI y su entorno.
 \item Endpoints básicos \texttt{POST /predict} y \texttt{GET /}; los definitivos figuran en el Anexo~\ref{cap:AnexoM}.
 \item Desarrollo de un servicio \emph{mock} para la ejecución de pruebas en entornos CI y el desarrollo de los otros componentes, sin necesidad de GPU.
 \item Dockerización del servicio con imagen base optimizada para PyTorch.
\end{itemize}

\subsection{Épica 4: Backend}

Tareas completadas:
\begin{itemize}[noitemsep]
 \item Preparación del entorno y las diferentes herramientas (\emph{linting}, análisis estático)
 \item Sistema de autenticación con registro, confirmación por correo electrónico e inicio de sesión con un token Bearer.
 \item Arquitectura CQRS/Event Sourcing con Commanded y PostgreSQL como \emph{event store}.
 \item Canal WebSocket para comunicación en tiempo real con el frontend.
 \item Pipeline de análisis para la recepción del informe médico, la llamada al AI Engine y el almacenamiento de los resultados.
 \item Soporte de internacionalización (i18n) en cinco idiomas (español, inglés, francés, italiano y alemán).
 \item Implementación de la suite de pruebas con ExUnit así como un informe de cobertura.
 \item Dockerización del entorno de desarrollo, pruebas y despliegue.
\end{itemize}

\subsection{Épica 5: Frontend}

Tareas completadas:
\begin{itemize}[noitemsep]
 \item Página de autenticación con formularios de registro e inicio de sesión.
 \item Interfaz de chat con renderizado de tarjetas tipadas (resumen y códigos).
 \item Panel lateral con historial de los diferentes análisis, así como de controles de sesión.
 \item Componentes de validación y rechazo de los códigos predichos, así como los motivos.
 \item Selector de idioma y carga de traducciones con el backend.
 \item Configuración centralizada de URLs mediante variables de entorno.
 \item Dockerización del entorno de desarrollo, pruebas y despliegue.
\end{itemize}

\subsection{Épica 6: Infraestructura y calidad}

Tareas completadas:
\begin{itemize}[noitemsep]
 \item Configuración de Docker Compose para todos los servicios.
 \item Makefile con targets para todo el ciclo de vida del proyecto.
 \item Configuración de commitlint y Husky para validación de commits.
 \item Workflow de GitHub Actions para CI.
 \item Dockerización del entorno de entrenamiento con soporte CPU y GPU.
 \item Implementación del entorno de despliegue con Docker Compose, Traefik y túnel de Cloudflare.
 \item Monitorización del servicio mediante Prometheus y Grafana.
 \item Gestión de errores utilizando el SaaS Sentry.
\end{itemize}

El listado completo de tareas del tablero, con sus identificadores y épica asociada, se recoge en el Anexo~\ref{cap:AnexoJ}.

\subsection{Orden de ejecución de las épicas}

Las épicas no se desarrollan de forma estrictamente secuencial. Los primeros pasos en todas las épicas serán siempre el establecimiento del entorno de desarrollo.

La adquisición de datos se aborda en segundo lugar, ya que el resto del sistema depende de ella: sin los CSV del corpus no hay entrenamiento posible. A continuación se trabaja en paralelo sobre el motor de IA y el backend, usando el \emph{mock} del AI Engine, para desacoplar ambas líneas de trabajo y poder avanzar en la integración del backend sin necesidad de disponer del modelo entrenado. El frontend se desarrolla progresivamente a medida que el backend expone nuevos endpoints y canales. La épica de modelo de clasificación es la de mayor duración y la que más variaciones presenta durante el desarrollo, dado que el rendimiento en validación determina si es necesario ajustar la arquitectura, los hiperparámetros o el preprocesamiento de los datos.

\subsection{Estimación de esfuerzo por épica}

La tabla~\ref{tab:planificacion} recoge la estimación de esfuerzo por épica expresada en puntos de historia relativos, donde cada punto equivale aproximadamente a media jornada de trabajo efectivo. Esta unidad relativa resulta más adecuada que las horas absolutas en un proyecto con dedicación variable e irregular, ya que refleja la complejidad percibida de cada épica con independencia de cuándo se realiza.

\begin{table}[h]
\centering
\begin{tabular}{lcc}
\hline
\textbf{Épica} & \textbf{Puntos estimados} & \textbf{Puntos reales} \\
\hline
Adquisición y preparación de datos & 8 & 8 \\
Modelo de clasificación & 16 & 24 \\
Servicio de inferencia & 6 & 10 \\
Backend & 16 & 16 \\
Frontend & 12 & 10 \\
Infraestructura y calidad & 4 & 4 \\
Redacción de la memoria & 20 & 22 \\
\hline
\textbf{Total} & \textbf{82} & \textbf{94} \\
\hline
\end{tabular}
\caption{Estimación y esfuerzo real por épica en puntos de historia.}
\label{tab:planificacion}
\end{table}

La épica de modelo de clasificación es la que más desviación ha sufrido. La diferencia de ocho puntos se atribuye a que se exploraron más combinaciones de \emph{batch size}, tasa de aprendizaje, función de pérdida y descongelación de las previstas (el proceso se documenta en el Anexo~\ref{cap:AnexoF}). Esta atribución es una estimación del autor y no una medición, como lo es la de los cuatro puntos de desviación del servicio de inferencia, que se achacan a adaptar el formato de entrada y salida del servicio a las configuraciones sucesivas del modelo. Este tipo de desviación es habitual en proyectos con componentes de investigación aplicada, donde el resultado de un experimento condiciona las decisiones que siguen.

\section{Costes y análisis de riesgos}

\subsection{Estimación de costes}

Los costes del proyecto se dividen en dos categorías: costes de personal y costes de recursos materiales. Dado que se trata de un TFG desarrollado por un único alumno, los costes de personal son los más significativos y los que mejor reflejan el esfuerzo real invertido.

\textbf{Costes de personal.} A partir de los 94 puntos de historia finalmente registrados (frente a los 82 estimados), asumiendo que cada punto equivale a media jornada de trabajo efectivo de cuatro horas, el volumen total de trabajo se sitúa en torno a las 376 horas. Se parte de un perfil sénior, con un salario bruto anual de entre 36.000 y 45.000~€ (supuesto del autor; la ganancia media anual de todos los trabajadores en España en 2024 fue de 29.540~€~\cite{ineeaes2024}, de modo que un perfil sénior de desarrollo queda por encima), unas cotizaciones a cargo de la empresa de en torno al 32\,\% (el 30,65\,\% de los tipos de contingencias comunes, desempleo, FOGASA, formación profesional y MEI de 2026~\cite{ordencotizacion2026}, más una prima de accidentes de trabajo supuesta de en torno al 1,4\,\%) y una jornada de 1.680 horas anuales (supuesto: las 40 horas semanales y los 30 días naturales de vacaciones del Estatuto de los Trabajadores~\cite{estatutotrabajadores2015} dejan unas 1.900 horas, de las que se descuentan festivos y ausencias). El coste horario es de entre unos 28 y 35 € con todos los costes sociales incluidos, y el coste de personal asciende a entre 10.528 y 13.160 €. Este rango representa el presupuesto de referencia si el proyecto se encargara a un profesional externo.

\textbf{Costes de recursos materiales.}
\begin{itemize}
 \item \textbf{Hardware de desarrollo (equipo y GPU)}: El desarrollo se ha realizado sobre hardware propio ya adquirido, por lo que no genera un coste adicional directo. No obstante, se imputa su amortización con dos supuestos del autor: un valor inicial de 2.500~€ para el equipo completo (AMD Ryzen 9 7900X y NVIDIA RTX 4080 incluidas) y una vida útil de cuatro años (625~€/año, unos 52~€ al mes). La amortización se imputa por el calendario del proyecto, de finales de enero a diciembre de 2026 (unos once meses), lo que da unos 573~€. La GPU forma parte de ese precio (1.469~€ de los 2.500~€, el precio de lanzamiento en España que publicó el fabricante~\cite{nvidiartx40anuncio2022}), por lo que su amortización no se cuenta aparte: de los 573~€, unos 337~€ corresponden a la GPU y unos 236~€ al resto del equipo. Repartidos entre las 376~horas de trabajo, el equipo cuesta en media unos 0,6~€ por hora de desarrollo (sin la GPU). Como referencia de mercado, alquilar una RTX 4090, de mayor potencia, cuesta desde 0,34~dólares por hora en RunPod~\cite{runpodprecios2026}, de modo que unas 95~horas de GPU (según se detalla más abajo) costarían desde unos 32~dólares, muy por debajo de los 337~€ imputados a la GPU. La diferencia se debe a que esa amortización se calcula por calendario y la GPU también sirve al sistema en funcionamiento, no solo al entrenamiento.
 \item \textbf{Electricidad del equipo}: El equipo funciona con energía solar de autoconsumo, de modo que el coste eléctrico directo es nulo. Como referencia a precio de la red convencional (0,14~€/kWh, tarifa del autor sin impuestos, que con el impuesto especial sobre la electricidad del 5,11\,\%~\cite{ley38impuestosespeciales1992} y el IVA del 21\,\%~\cite{ley37iva1992} resulta en unos 0,18~€/kWh) y con consumos supuestos por el autor, la electricidad del proyecto sumaría unos 143~€ (800~kWh), en tres partidas: \emph{(a)}~unas 95~horas de GPU a plena carga (450~W, 43~kWh; la GPU consume hasta 320~W~\cite{nvidiartx4080} y la CPU hasta 170~W~\cite{amdryzen7900x}), unos 8~€; \emph{(b)}~las 376~horas de trabajo de desarrollo (150~W, 56~kWh), unos 10~€; y \emph{(c)}~el servicio encendido de forma continua desde mayo hasta diciembre de 2026, ocho meses (unas 5.840~horas a 120~W con el modelo cargado, 701~kWh), unos 125~€. Las horas de GPU suman tres partidas de distinta fiabilidad: \emph{(i)}~38,8~horas de ajuste fino, registradas en \texttt{training\_runs.csv} (suma de \texttt{total\_seconds} de las 78 ejecuciones); \emph{(ii)}~unas 37~horas de preentrenamiento débil, que no se registran, estimadas a partir de las 550~épocas de preentrenamiento de las 56 ejecuciones que lo usaron (columna \texttt{pretrain\_epochs}) a unos 4~minutos por época, tiempo extrapolado de una medición con tres ejecuciones simultáneas (unas 2~horas por 10~épocas, a 0,6--0,9~s por paso frente a unos 0,25~s en solitario); y \emph{(iii)}~unas 20~horas, estimadas por el autor, de los cuadernos de la fase de selección de modelos (\texttt{training/bert-classifier/*.ipynb}), cuyo tiempo no se registró. Las partidas (ii) y (iii) son estimaciones y la suma no incluye las evaluaciones y mediciones sueltas ni la carga de los modelos.
 \item \textbf{Software}: Todas las herramientas, \emph{frameworks} y modelos utilizados son de código abierto y gratuitos. El acceso al modelo RigoBERTa desde HuggingFace Hub requiere un token de usuario, cuya obtención es completamente gratuita.
 \item \textbf{Infraestructura de red y almacenamiento}: El corpus CodiEsp y los artefactos del modelo se almacenan localmente. Por tanto, no se han generado costes de almacenamiento en la nube ni de ancho de banda significativos.
\end{itemize}

La Tabla~\ref{tab:costes} resume la estimación total, con el coste de personal como partida dominante.

\begin{table}[h]
\centering
\begin{tabular}{lrr}
\hline
\textbf{Concepto} & \textbf{Mínimo} & \textbf{Máximo} \\
\hline
Personal (376\,h × 28--35\,€/h) & 10.528\,€ & 13.160\,€ \\
Amortización del equipo, GPU incluida (11 meses) & 573\,€ & 573\,€ \\
Electricidad del equipo (energía solar; equivalente en red) & 0\,€ & 143\,€ \\
Software e infraestructura & 0\,€ & 0\,€ \\
\hline
\textbf{Total} & \textbf{11.101\,€} & \textbf{13.876\,€} \\
\hline
\end{tabular}
\caption{Resumen de costes estimados del proyecto.}
\label{tab:costes}
\end{table}

El coste total estimado del proyecto, incluyendo personal y materiales, se sitúa entre \textbf{11.100 y 13.900~€}, siendo el coste de personal el componente dominante, con alrededor del 95\,\% del total.

\subsection{Análisis de riesgos}

La identificación y gestión de riesgos es una parte esencial del plan de gestión de cualquier proyecto software, más aún en proyectos con componentes de investigación aplicada, donde el resultado de los experimentos no está garantizado de antemano. A continuación se enumeran los riesgos identificados durante la planificación del proyecto, clasificados por su probabilidad de ocurrencia (Baja, Media o Alta) y por el impacto que tendrían sobre el desarrollo si se materializaran (Bajo, Medio o Alto).

La tabla~\ref{tab:riesgos} recoge los principales riesgos identificados, su probabilidad de ocurrencia, el impacto estimado y las medidas de mitigación adoptadas.

\begin{table}[h]
\centering
\small
\begin{tabular}{P{5cm}ccP{5.2cm}}
\hline
\textbf{Riesgo} & \textbf{Prob.} & \textbf{Impacto} & \textbf{Mitigación} \\
\hline
Rendimiento insuficiente del modelo sobre CodiEsp & Media & Alto & Evaluación temprana sobre el conjunto de validación; posibilidad de cambiar el modelo base o ajustar hiperparámetros. \\
\hline
Conjuntos de datos con muy pocos ejemplos para ciertos códigos (desequilibrio de clases) & Alta & Medio & Análisis previo de la distribución; ponderación de clases en la función de pérdida. \\
Coste computacional excesivo de entrenamiento sin GPU & Alta & Medio & Configuración Docker con soporte GPU; los entornos en la nube serían una alternativa que no se ha probado. \\
\hline
Cambios en la API del portal eCIE-maps que rompan el \emph{scraping} & Baja & Medio & Los CSV ya generados se almacenan en el repositorio; el \emph{scraping} es una fase de adquisición puntual, no continua. \\
\hline
Indisponibilidad del corpus CodiEsp & Baja & Alto & Copia local de los datos incluida en el repositorio de entrenamiento. \\
\hline
\end{tabular}
\caption{Análisis de riesgos del proyecto.}
\label{tab:riesgos}
\end{table}

De los riesgos identificados, los dos que mayor atención requieren durante el desarrollo son el rendimiento del modelo y el desequilibrio de clases. El primero porque condiciona directamente la utilidad del sistema. El criterio de éxito no era superar a los mejores sistemas publicados, sino situarse dentro del rango de resultados de la comparativa CodiEsp-D (Tabla~\ref{tab:codiesp-comparativa}) y declarar con honestidad en qué puesto queda. El segundo porque es una característica inherente a cualquier corpus de diagnósticos clínicos reales, donde la distribución de frecuencias sigue una ley de potencia: un pequeño conjunto de códigos de enfermedades comunes acapara la mayoría de los ejemplos de entrenamiento, mientras que los códigos de enfermedades raras o procedimientos poco frecuentes cuentan con muy pocos ejemplos. Este desequilibrio puede llevar a que el modelo aprenda a predecir siempre los códigos más frecuentes y a ignorar sistemáticamente los menos comunes, lo que penaliza las métricas macro y degrada la utilidad clínica del sistema para los casos atípicos. La ponderación de clases en la función de pérdida y el análisis previo de la distribución del corpus son las medidas de mitigación más directas para este riesgo.
